\documentclass{article}

\usepackage[preprint]{neurips_2026}

\usepackage[utf8]{inputenc}
\usepackage[T1]{fontenc}
\usepackage{hyperref}
\usepackage{url}
\usepackage{booktabs}
\usepackage{amsmath}
\usepackage{amssymb}
\usepackage{amsfonts}
\usepackage{graphicx}
\usepackage{nicefrac}
\usepackage{microtype}
\usepackage{xcolor}
\usepackage{multirow}
\usepackage{algorithm}
\usepackage{algpseudocode}

\title{Dependency-Coherent Coordination: Flow-Based Partitioning Improves
    Multi-Agent Traffic Signal Control Across All Coordination Levels}

\author{%
  [Author(s) — anonymised for review]
}

\begin{document}

\maketitle

% ============================================================
\begin{abstract}
% TODO — write after all sections are drafted.
% Target: 150–200 words.
% Structure: (1) problem + gap, (2) method, (3) key result, (4) significance.
%
% Draft placeholder:
Coordinating traffic signal agents in multi-agent reinforcement learning
(MARL) requires partitioning intersections into coordination groups, yet
existing partitions use road topology rather than vehicle flow structure.
We ask: \emph{which inter-agent dependencies does a partition preserve or
    sever, and does preserving more dependencies improve coordination?}
We apply InfoMap — an information-theoretic community detection algorithm
grounded in origin-destination (OD) vehicle flow — to define
flow-coherent coordination modules, and evaluate its effect systematically
across three coordination coupling levels: reward shaping (IPPO), centralised
critic (MAPPO), and graph attention (CoLight).
We introduce the \emph{dependency retention ratio}~$\Psi$, a partition quality
metric computed from mutual information and Jacobian-based dependency
instruments, and show that InfoMap achieves higher $\Psi$ than topology-based
partitions (METIS) across all three levels.
Experiments on real SUMO scenarios demonstrate up to $2.85\times$ reduction in
average vehicle waiting time, with verified spatial MI advantage (H1) and
spill-back containment (H2) as the causal mechanism.
\end{abstract}

% ============================================================
\section{Introduction}
\label{sec:intro}
% Para 1 — Problem: TSC at scale; coordination graph design as the open question
% Para 2 — Gap (two layers):
%   Structural: existing partitions (METIS, Laplacian, feudal) use topology, not OD flow
%   Theoretical: dynamic graph attention learns structure implicitly but provides no
%                principled account of which dependencies are preserved or severed
% Para 3 — Insight + method: OD flow is the natural basis; InfoMap is the
%           information-theoretic criterion; coordination coupling spectrum (L1/L2/L3)
% Para 4 — Contributions (bulleted, 4 items)
\textbf{[To be written — Step 6]}

% ============================================================
\section{Background}
\label{sec:background}

\subsection{Network-Distributed POMDP and the Coordination Graph}
\label{sec:ndpomdp}

Each signalised intersection $i \in \mathcal{V}$ observes a local state
$o_i$ (queue lengths, current phase, elapsed phase time) and selects a
phase action $a_i \in \mathcal{A}_i$.
We model the joint problem as a \emph{network-distributed POMDP} (ND-POMDP),
in which a coordination graph $\mathcal{G}_C = (\mathcal{V}, \mathcal{E}_C)$
encodes which agent pairs share explicit interaction terms in the joint reward.
Three structural properties of TSC motivate this formulation.

\textbf{(a) Reward locality with spill-back coupling.}
The immediate reward $r_i$ depends primarily on queues and delays at
intersection $i$, making it nearly local.
However, when downstream lanes saturate, backward-propagating queue waves
couple $r_i$ to the actions at neighbouring intersection $j$ via
spill-back~\citep{dependencydyn2025}.
The ND-POMDP captures this through interaction terms:
\begin{equation}
    R(\mathbf{a})
    \;=\; \sum_{i \in \mathcal{V}} r_i(o_i, a_i)
    \;+\; \sum_{(i,j)\in\mathcal{E}_C} r_{ij}(o_i, o_j, a_i, a_j).
    \label{eq:ndpomdp}
\end{equation}
Edges in $\mathcal{G}_C$ mark pairs where the cross-agent term $r_{ij}$ is
non-negligible; omitting a spill-back edge misattributes joint outcomes to
independent agents.

\textbf{(b) Tractability via maximum clique size.}
Exact coordination over a clique of $\omega$ agents requires evaluating
$|\mathcal{A}|^\omega$ joint actions.
Partitioning $\mathcal{V}$ into modules of bounded size
$|M_k| \leq \omega_{\max}$ keeps this tractable; in TSC networks
$\omega_{\max} = 4$--$6$ is sufficient in practice.

\textbf{(c) $\mathcal{G}_C$ is a free design variable.}
Unlike environments where the coordination graph is determined by the physics,
in TSC the coordination modules are explicitly constructed before training.
This makes partition selection the central algorithmic question: which graph
minimises the coordination cost severed by module boundaries?
Attribute~(a) answers what ``coordination cost'' means: the mass of spill-back
interaction terms $r_{ij}$ cut by the partition.
Attributes~(a)–(c) together motivate using OD vehicle flow as the partition
criterion, and Section~\ref{sec:infomap} makes this precise.

\subsection{The Map Equation}
\label{sec:mapeq}

InfoMap~\citep{rosvall2008maps,rosvall2009mapequation} finds a partition by
minimising the description length of a random walk on a weighted directed graph:
\begin{equation}
    L(\mathcal{M})
    \;=\; q_{\circlearrowleft} H(\mathcal{Q})
    \;+\; \sum_{k=1}^{K} p_k^{\circlearrowleft} H(\mathcal{P}_k),
    \label{eq:mapeq}
\end{equation}
where $q_{\circlearrowleft}$ is the aggregate rate of inter-module random-walk
transitions, $H(\mathcal{Q})$ is the entropy of the between-module codebook,
$p_k^{\circlearrowleft}$ is the within-module random-walk rate for $M_k$, and
$H(\mathcal{P}_k)$ is the within-module codebook entropy.
Minimising $L(\mathcal{M})$ minimises the average bits required to describe
the walker's trajectory, which is achieved by grouping nodes that exchange the
most flow into the same module --- so that the walker rarely crosses module
boundaries and the between-module term $q_{\circlearrowleft}H(\mathcal{Q})$
stays small.

Applied to $\mathcal{G}_{\mathrm{OD}}$, InfoMap groups intersections that
exchange the most vehicle volume.
As we show in Section~\ref{sec:infomap}, these are exactly the pairs most
likely to be coordination-necessary via spill-back, connecting the
information-theoretic objective of Eq.~\eqref{eq:mapeq} to the traffic-theoretic
requirement of Eq.~\eqref{eq:ndpomdp}.
Unlike modularity maximisation~\citep{newmangirvan2004}, the map equation has
no resolution limit~\citep{edler2023community}: it detects modules at any
scale present in the data without merging fine-grained clusters that are
genuinely distinct.

% ============================================================
\section{Related Work}
\label{sec:related}
% Para 1 — Network partitioning for TSC:
%   laplacian2009ieee, odflow2019partition, mfdpartition2020,
%   communitydetection2021plosone, ma2022feudal, liu2023gplight
%   → structural/geometric criteria; no OD flow objective; no dependency cost analysis
% Para 2 — MARL communication topology (compressed):
%   jiang2018atoc, i2c2021, hetib2025, wei2019colight, liu2023gplight,
%   xu2021hilight, zhu2026halo, expocomm2025, li2025topology
%   → learned or fixed; dynamic graph attention implicit; no dependency accounting
% Para 3 — GNN dependency + MI in MARL (compressed):
%   lrd2025, abnar2020rollout, stmarl2020, dependencydyn2025, wang2020imac
%   → Ψ, D^MI, D^Q, D^G adapt these frameworks to partition quality evaluation
\textbf{[To be written — Step 5]}

% ============================================================
\section{Flow-Coherent Dependency Partitioning}
\label{sec:method}

\subsection{OD Flow Graph Construction}
\label{sec:odgraph}

We model the traffic network as a directed weighted graph
$\mathcal{G}_{\mathrm{OD}} = (\mathcal{V}, \mathcal{E}, \mathbf{W})$,
where $\mathcal{V}$ is the set of $N$ signalised intersections and $W_{ij} \geq 0$
is the observed vehicle count travelling from intersection $i$ to intersection $j$
over a representative time window.
We construct $\mathbf{W}$ from simulation trip logs: for each completed vehicle
trip, the sequence of intersections visited is recorded and the corresponding
directed edge weights incremented.
A warm-up period of approximately one hour under a uniformly random policy provides
a flow snapshot that reflects network demand structure prior to any learned behaviour.

The key inductive bias is that $\mathbf{W}$ captures \emph{demand} rather than
\emph{road topology}: two intersections connected by a high-volume corridor appear
as a strong edge regardless of their hop distance on the road graph.
This is what separates OD-flow partitioning from METIS on the adjacency matrix
and from Laplacian spectral methods~\citep{laplacian2009ieee}: both of the latter
assign edge weight by physical connectivity, not by how many vehicles actually
travel between node pairs.

\subsection{InfoMap Partitioning and the Dependency Gate}
\label{sec:infomap}

Given $\mathcal{G}_{\mathrm{OD}}$, we apply the Infomap
algorithm~\citep{rosvall2008maps,rosvall2009mapequation} to obtain a partition
$\mathcal{M} = \{M_1, \ldots, M_K\}$ by minimising the map equation:
\begin{equation}
    L(\mathcal{M})
    = q_{\circlearrowleft} H(\mathcal{Q})
    + \sum_{k=1}^{K} p_k^{\circlearrowleft} H(\mathcal{P}_k)
    \label{eq:mapequation}
\end{equation}
where $q_{\circlearrowleft}$ is the rate of inter-module random-walk transitions
and $H(\mathcal{Q})$ is the entropy of the between-module codebook.
Minimising $L(\mathcal{M})$ maximises the fraction of vehicle-flow steps retained
within modules — equivalent, in the ND-POMDP interpretation, to minimising
cross-module coordination overhead while maximising within-module coordination
coherence.
Crucially, Infomap has no resolution limit~\citep{edler2023community}: it recovers
modules at any scale present in the data, unlike modularity
maximisation~\citep{newmangirvan2004}.

We verify two empirical conditions on the partition before MARL training.
Results are reported as supporting evidence; failure on a given network
indicates a boundary condition requiring separate investigation.

\medskip
\noindent\textbf{H1 (Spatial MI Advantage).}
Within-module agent pairs share higher mutual information about future traffic
states than cross-module pairs:
\begin{equation}
    \mathbb{E}_{i,j\in M_k}\!\left[I(X_i(t);\, X_j(t{+}\tau))\right]
    >
    \mathbb{E}_{i\in M_k,\, l\notin M_k}\!\left[I(X_i(t);\, X_l(t{+}\tau))\right]
    \label{eq:h1}
\end{equation}
averaged over propagation lags $\tau \in \{1, 5, 15, 30\}$ steps.
Verified via Wilcoxon signed-rank test ($p < 0.05$) on KSG-estimated
MI~\citep{kraskov2004ksg}; visualised as a $|\mathcal{V}|\times|\mathcal{V}|$
heatmap with InfoMap module boundaries overlaid.

\medskip
\noindent\textbf{H2 (Spill-back Containment).}
Spill-back congestion is more likely to remain within InfoMap modules than to
cross module boundaries:
\begin{equation}
    \Pr\!\left(\text{spill-back}_{i\to j}
    \;\middle|\; i,j \in M_k\right)
    >
    \Pr\!\left(\text{spill-back}_{i\to l}
    \;\middle|\; i \in M_k,\; l \notin M_k\right).
    \label{eq:h2}
\end{equation}
\citet{dependencydyn2025} establish that inter-agent coordination is necessary
\emph{if and only if} spill-back connects two agents' queues.
Verified via a one-sample proportion test on SUMO queue trajectory data.
The apparent directional asymmetry --- OD flow travels $i \to j$ (forward) while
spill-back propagates $j \to i$ (backward) --- does not invalidate H2: a
high-volume corridor $W_{ij}$ drives high arrival rates at $j$ which, under
near-saturated demand, trigger backward-propagating queue waves toward
$i$~\citep{dependencydyn2025}.
The coordination-necessary pair and the high-OD-flow pair are thus the same pair
on the same corridor.

\medskip
H2 is the traffic-theoretic \emph{mechanism} behind H1: spill-back containment
causes within-module states to co-vary via shared congestion propagation, raising
$I(X_i;\,X_j)$ above the cross-module baseline.
Together, H2$\,\Rightarrow\,$H1 imply the main hypothesis~H: an InfoMap partition
retains more inter-agent spatial dependency within modules than a topology-based
partition of equal size --- formally quantified by $\Psi$ (Section~\ref{sec:psi}).

\subsection{Spatial Dependency Instruments and $\Psi$}
\label{sec:psi}

To measure how much dependency a partition retains, we need per-pair dependency
weights $D_{ij}$ that (i)~are computable without presupposing the partition,
(ii)~yield a scalar per pair $(i,j)$, and (iii)~can be aggregated into a
partition-level score. We introduce three complementary instruments, each matched
to the coordination mechanism at a given coupling level.

\paragraph{Physical dependency (model-free).}
H1 directly supplies a model-free dependency weight via the mutual information
between agent state observations:
\begin{equation}
    D^{\mathrm{MI}}_{ij}
    = \hat{I}\!\left(X_i(t);\, X_j(t{+}\tau)\right)
    \label{eq:dmi}
\end{equation}
estimated by the KSG $k$-nearest-neighbour estimator~\citep{kraskov2004ksg}
from warm-up trajectories.
$D^{\mathrm{MI}}_{ij}$ is model-free and applicable at L1: at the reward
coupling level (IPPO/IDQN), agents act independently and there is no
cross-agent input path in the policy network, making Jacobian-based measures
structurally inapplicable.
Spatial MI dominance in TSC networks is established by~\citet{stmarl2020},
and its mechanistic link to spill-back is established by~\citet{dependencydyn2025},
grounding $D^{\mathrm{MI}}$ as the natural physical dependency measure.

\paragraph{Learned dependency (Jacobian on the global model).}
For models that receive cross-agent inputs, we adapt the Jacobian influence
framework of~\citet{lrd2025} to the MARL dependency measurement setting.
For a model $F$ trained \emph{without} partition constraints — the global model
— define the spatial influence of agent~$j$ on agent~$i$:
\begin{equation}
    D_{i\leftarrow j}(F)
    = \sum_{a,k}
    \left|\frac{\partial F_{i,a}}{\partial X_{j,k}}\right|
    \label{eq:djac}
\end{equation}
where $X_{j,k}$ is feature~$k$ of agent~$j$'s observation and $F_{i,a}$ is
output dimension~$a$ at agent~$i$.
The \emph{global-model restriction} is essential: in a partitioned model,
cross-module inputs are absent by construction, making cross-module Jacobians
identically zero — a structural artefact, not evidence of low dependency.
We instantiate Eq.~\eqref{eq:djac} at two levels:
\begin{align}
    D^{Q}_{i\leftarrow j}
     & = D_{i\leftarrow j}(V^{\mathrm{global}})
     &                                                         & \text{(global MAPPO critic, L2)}
    \label{eq:dq}                                                                                          \\
    D^{G}_{i\leftarrow j}
     & = D_{i\leftarrow j}(\mathbf{h}^{(L),\mathrm{global}}_i)
     &                                                         & \text{(global CoLight GAT embedding, L3)}
    \label{eq:dg}
\end{align}
Whereas \citet{lrd2025} use Jacobian profiles to ask whether a GNN suffers
from over-squashing, we use the same instrument to ask whether a proposed partition
aligns with what the model has learned to depend on — a distinct scientific question.

\paragraph{Hop-distance profile.}
For either instrument, aggregate dependency by road-graph hop distance $\rho(i,j)$:
\begin{equation}
    \bar{D}_h
    = \frac{1}{|\mathcal{V}|}\sum_{i}\sum_{j:\,\rho(i,j)=h} D_{i\leftarrow j}
    \label{eq:barDh}
\end{equation}
This is the MARL analogue of the per-hop influence profile $\bar{T}_h$
of~\citet{lrd2025}: $\bar{D}_h$ quantifies the average spatial dependency
from agents exactly $h$ road-graph hops away, and its decay rate determines
the minimum module radius needed for effective coordination.

\paragraph{Dependency retention ratio $\Psi$.}
Given dependency matrix $\mathbf{D}$ and partition $\mathcal{P}=\{P_1,\ldots,P_K\}$:
\begin{equation}
    \Psi(\mathcal{P},\,\mathbf{D})
    = \frac{\displaystyle\sum_{k}\sum_{i,j\,\in\, P_k} D_{ij}}
    {\displaystyle\sum_{i,j} D_{ij}}
    \label{eq:psi}
\end{equation}
$\Psi \in [0,1]$ is the \emph{dependency retention ratio}: the fraction of total
inter-agent dependency captured within module boundaries.
High $\Psi$ means the partition cuts predominantly low-dependency edges —
preserving coordination capacity — while low $\Psi$ means high-dependency edges
are severed, imposing coordination cost.

\paragraph{Relation to existing partition quality metrics.}
Graph modularity $Q$~\citep{newmangirvan2004} and normalised
cut~\citep{shimalik2000} are the standard partition quality metrics.
$\Psi$ differs from both in two respects.
First, $\Psi$ requires no null model: $\mathbf{D}$ is already a semantically
meaningful quantity (MI or Jacobian sensitivity), not raw structural adjacency
for which a random-graph null is appropriate.
Second, $\mathbf{D}$ is asymmetric and non-negative (direction matters: agent $j$
influences $i$ via traffic flow), whereas $Q$ and normalised cut assume undirected
graphs.
By computing $\Psi(\mathcal{P}_{\mathrm{InfoMap}},\mathbf{D})$ vs.\
$\Psi(\mathcal{P}_{\mathrm{METIS}},\mathbf{D})$ on the \emph{same} $\mathbf{D}$,
we obtain a direct, architecture-independent measure of which partition better
preserves the inter-agent dependency structure.

\subsection{Coordination Coupling Spectrum}
\label{sec:levels}

We evaluate the InfoMap partition across three coordination coupling levels,
using established MARL algorithms as probes. The partition is held fixed across
all three; only the coordination mechanism varies.
This design isolates the partition's contribution from algorithm-specific effects,
and tests whether the dependency retention advantage of InfoMap is
architecture-agnostic.

\begin{table}[t]
    \centering
    \caption{Coordination coupling spectrum. The same InfoMap partition is applied
        at each level. The dependency instrument column indicates which metric from
        Section~\ref{sec:psi} measures partition quality at that level, and which
        operational prediction (OP) links $\Psi$ to expected task improvement.}
    \label{tab:levels}
    \small
    \begin{tabular}{llllll}
        \toprule
        Level & Coupling      & Mechanism                   & Algorithm & Instrument        & Prediction \\
        \midrule
        L1    & Implicit      & Regional reward averaging   & IPPO/IDQN & $D^{\mathrm{MI}}$ & OP1        \\
        L2    & Semi-explicit & Regional centralised critic & MAPPO     & $D^{Q}$           & OP2        \\
        L3    & Explicit      & Regional graph attention    & CoLight   & $D^{G}$           & OP3        \\
        \bottomrule
    \end{tabular}
\end{table}

\noindent\textbf{Operational predictions.}
The following are the primary empirical targets of our experimental design,
derived from H1 and H2 via the causal chain H2$\,\Rightarrow\,$H1$\,\Rightarrow\,$H:
\begin{description}
    \item[\textbf{OP1 (L1):}]
          $\Psi(\mathcal{P}_{\mathrm{InfoMap}}, D^{\mathrm{MI}}) >
              \Psi(\mathcal{P}_{\mathrm{METIS}}, D^{\mathrm{MI}})$.
          Within-module reward averaging operates on correlated signals, reducing
          gradient noise and training instability.
    \item[\textbf{OP2 (L2):}]
          $\Psi(\mathcal{P}_{\mathrm{InfoMap}}, D^{Q}) >
              \Psi(\mathcal{P}_{\mathrm{METIS}}, D^{Q})$.
          The regional critic captures dominant value dependencies without modelling
          irrelevant cross-module states, reducing credit-assignment variance.
    \item[\textbf{OP3 (L3):}]
          $\Psi(\mathcal{P}_{\mathrm{InfoMap}}, D^{G}) >
              \Psi(\mathcal{P}_{\mathrm{METIS}}, D^{G})$, observed via a lower
          influence-weighted attention range $\hat{\rho}_u$ for regional vs.\ global
          CoLight (Section~\ref{sec:attnrange}), provided module size is
          sufficient (${\geq}4$ agents per module).
\end{description}

\noindent The three OPs together constitute \textbf{Hypothesis H}: an InfoMap
partition achieves higher $\Psi$ than METIS across all three dependency instruments,
and this advantage translates to better task performance at each coupling level.

% ============================================================
\section{Experiments}
\label{sec:experiments}
% §5.1 Setup: Cologne8 (8 int.), Ingolstadt21 (21 int.), [6×6 grid if available]
%      Baselines: L0 (no-partition) → L1.5 (random-sparse) → METIS → InfoMap
%      Protocol: rliable — IQM ± stratified bootstrap CI, ≥5 seeds, AUC-LC + ATT
% §5.2 H1/H2 verification — Fig 2: MI heatmap + module boundaries (run before training)
% §5.3 Main results — Table 1 + Fig 3 learning curves
%      IQM ± CI: InfoMap / METIS / None × L1 / L2 / L3, primary: Ingolstadt21
% §5.4 Ψ predicts performance — Fig 4: scatter ΔΨ vs ΔAUC-LC
% §5.5 Attention range analysis — Fig 5: D̄_h profile + rollout heatmaps
\label{sec:attnrange}
% §5.6 Ablations: module size k, L1.5 sparsification control, reward annealing α(e)
\textbf{[To be written — Step 7]}

% ============================================================
\section{Conclusion}
\label{sec:conclusion}
% Mechanism: InfoMap → high Ψ → better L1/L2/L3
% Verified causal chain: H2 ⇒ H1 ⇒ H ⇒ OP1/OP2/OP3
% Boundary: L3 requires ≥4 agents/module
% Limitations: warm-up OD approximation; static partition
% Future: dynamic repartitioning; real-world OD data
\textbf{[To be written — Step 8]}

% ============================================================
\begin{ack}
    [Omitted in anonymised submission.]
\end{ack}

% ============================================================
\section*{References}
{
    \small
    \bibliographystyle{plainnat}
    \bibliography{refs}
}

% ============================================================
\appendix

\section{Algorithm: InfoMap Coordination Graph Construction}
\label{app:algorithm}

H1 and H2 are \emph{manipulation checks}, not engineering controls.
The algorithm always returns the InfoMap partition; the verification results are
reported as supporting evidence and logged for transparency.
A failure on a given network indicates a network-specific boundary condition
(e.g.\ very low overall OD volume, or grid topology with uniform flow) and is
reported separately.

\begin{algorithm}[H]
    \caption{Flow-coherent coordination graph construction with dependency verification}
    \label{alg:infomap}
    \begin{algorithmic}[1]
        \Require Traffic network $\mathcal{G} = (\mathcal{V}, \mathcal{E})$,
        number of modules $K$, warm-up steps $T_w$
        \Ensure  Partition $\mathcal{M} = \{M_1, \ldots, M_K\}$,\;
        dependency matrix $\mathbf{D}^{\mathrm{MI}}$,\;
        verification results $(p_{\mathrm{H1}},\, p_{\mathrm{H2}})$
        \State Run $T_w$ steps under a uniformly random policy; log all vehicle trips
        \State Construct $\mathbf{W}$: $\;W_{ij} \leftarrow$ vehicle count $i \to j$
        \State Apply Infomap to $\mathcal{G}_{\mathrm{OD}} = (\mathcal{V}, \mathcal{E}, \mathbf{W})$
        $\;\Rightarrow\;$ partition $\mathcal{M} = \{M_1,\ldots,M_K\}$
        \State Estimate $D^{\mathrm{MI}}_{ij}$ via KSG estimator~\citep{kraskov2004ksg}
        at $\tau \in \{1,5,15,30\}$ steps
        \Statex \hspace{1em}\textit{// H2 manipulation check}
        \State Extract spill-back events from SUMO queue trajectories
        \State Compute within-module and cross-module spill-back proportions;
        run one-sample proportion test $\;\Rightarrow\; p_{\mathrm{H2}}$
        \Statex \hspace{1em}\textit{// H1 manipulation check}
        \State Classify agent pairs as within-module or cross-module under $\mathcal{M}$
        \State Compare $D^{\mathrm{MI}}_{ij}$ distributions;
        run Wilcoxon signed-rank test $\;\Rightarrow\; p_{\mathrm{H1}}$
        \State \textbf{Report} $(p_{\mathrm{H1}},\, p_{\mathrm{H2}})$;
        generate MI heatmap with module boundaries (Fig.~2 in main text)
        \State \Return $\mathcal{M}$,\; $\mathbf{D}^{\mathrm{MI}}$,\;
        $(p_{\mathrm{H1}},\, p_{\mathrm{H2}})$
    \end{algorithmic}
\end{algorithm}

% ============================================================
\section{Regional MARL Architecture}
\label{app:arch}

All three coordination levels share a common regional infrastructure.
Partition assignments are stored in \texttt{cfg.management} (one list of
agent IDs per module), which is populated from the InfoMap or METIS output
before any training begins.
A central helper module (\texttt{agents/common/marl.py}) provides
\texttt{transform\_reward\_array(scope, rewards, cfg)} --- supporting
\texttt{local}, \texttt{global\_mean}, and \texttt{regional\_mean} scopes ---
and \texttt{CentralizedInputBuilder}, which assembles critic observations for
\texttt{global}, \texttt{regional}, or \texttt{local} scopes and zero-pads
uneven regional input sizes to a common dimension.

\subsection{L1 --- Regional Reward (IPPO)}
\label{app:arch:ippo}

Regional IPPO applies flow-partition structure at the reward level only.
Each agent maintains an independent local actor-critic PPO learner
(\texttt{agents/policy/pfrl\_ppo.py}); there is no cross-agent parameter
sharing or communication.
Before each learner's \texttt{observe()} call, individual scalar rewards
are passed through \texttt{transform\_reward\_array} with
\texttt{reward\_scope: regional\_mean}, replacing each agent's reward with
the arithmetic mean of rewards within its InfoMap module.
Unassigned agents (singletons not covered by \texttt{cfg.management}) retain
their individual reward.
The default configuration (\texttt{config/agent.yaml}) sets
\texttt{reward\_scope: local}; regional behaviour is enabled by
\texttt{reward\_scope: regional\_mean}.
IDQN (\texttt{agents/action\_value/pfrl\_dqn.py}) follows the identical
reward-rewriting path.

\paragraph{Mechanism.}
Regional averaging reduces credit-assignment noise when agents within the
same module experience correlated outcomes (high $D^{\mathrm{MI}}$).
It is harmful when modules group uncorrelated agents, because the averaged
signal is less informative than the individual reward.
This makes IPPO performance at L1 a direct diagnostic of partition quality.

\subsection{L2 --- Regional Critic (MAPPO)}
\label{app:arch:mappo}

MAPPO is implemented in \texttt{agents/policy/actor\_critic.py} on top of
\texttt{ActorCriticBase}.
Actors use local observations and mask invalid actions for heterogeneous
action spaces (shared actor-network weights across agents).
The critic input is assembled by \texttt{CentralizedInputBuilder} under
\texttt{critic\_scope: regional}: each agent's critic receives the concatenated
observations of all agents in its InfoMap module, zero-padded to the maximum
regional size.
Reward sharing is controlled separately via \texttt{reward\_scope}; the
standard regional MAPPO configuration uses both
\texttt{critic\_scope: regional} and \texttt{reward\_scope: regional\_mean}.

Global MAPPO (\texttt{critic\_scope: global}) feeds all $N$ agents'
observations to a single shared critic.
For $N = 21$ (Ingolstadt21), the global critic input dimension is $21 \times
    d_{\mathrm{obs}}$, making it substantially harder to train than a regional
critic of dimension $|M_k| \times d_{\mathrm{obs}}$ where $|M_k| \approx 4$--$5$.

\paragraph{Mechanism.}
The regional critic captures dominant value dependencies
($D^Q_{i\leftarrow j}$ high within module) without modelling irrelevant
cross-module states, reducing critic approximation error and policy-gradient
variance.
On networks where a global critic is impractical (large $N$), regional InfoMap
critics can \emph{outperform} global CTDE, as observed on Ingolstadt21.

\subsection{L3 --- Regional Attention (CoLight)}
\label{app:arch:colight}

CoLight is implemented in \texttt{agents/action\_value/colight.py}.
The topology is built from \texttt{signal.yaml} downstream links as
bidirectional graph edges.
Two distinct regional variants exist; they test different scientific questions.

\paragraph{Variant A: PartitionedAttentionStack.}
Enabled by \texttt{colight\_partitioned\_message\_passing: True}.
One GAT attention stack per module is created; all stacks share the same
underlying network weights and a single replay buffer.
Attention is structurally restricted to within-module edges --- agents in
different modules cannot attend to each other.
This is the correct variant for testing whether InfoMap topology improves
the quality of \emph{communication}, as the learned $D^G$ Jacobians are
computed on a globally-unified model but message-passing is constrained.

\paragraph{Variant B: CoLightRegional.}
Separate agent class that instantiates one independent
\texttt{\_CoLightRegionLearner} per partition.
Each regional learner has its own CoLight network, target network, optimizer,
replay buffer, action mask, regional edge index, and $\epsilon$-schedule.
\texttt{CoLightRegional.observe()} converts individual rewards to the mean
reward within each region before inserting transitions into that region's
replay buffer.
This variant tests whether divide-and-conquer training by region is beneficial;
it is \emph{not} equivalent to restricting communication in a global model.
When module size is small ($|M_k| \lesssim 3$), Variant B creates highly
under-sampled regional learners with unstable target networks.

\paragraph{Fallback.}
If \texttt{cfg.management} is absent, both variants fall back to one global
region covering all agents.

% ============================================================
\section{Preliminary Results (Work in Progress)}
\label{app:results}

All values are \textbf{final-episode eval metrics} (single seed) unless noted.
Episode counts below 1500 indicate runs that have not yet completed.
Lower is better for all metrics.
Delay $=$ timeLoss (s/veh); Trip $=$ duration (s/veh);
Wait $=$ waitingTime (s/veh); Queue $=$ mean queue length (veh/intersection).

\subsection{Cologne8 (C8, 8 intersections, k\,=\,3, seed\,0, 1500\,ep)}
\label{app:results:c8}

\begin{table}[H]
\centering
\caption{Cologne8 --- full 4-metric results. All runs complete (1500\,ep)
    except IDQN baseline (crashed ep\,1340, marked $\dagger$).
    Seed 7 used for CoLight runs; all others seed 0.}
\label{tab:c8_full}
\small
\begin{tabular}{llrrrr}
    \toprule
    Algorithm & Condition                        & Delay & Trip   & Wait  & Queue \\
    \midrule
    \multirow{3}{*}{IPPO}
              & No partition (baseline)          & 39.17 & 104.35 & 19.32 & 2.22  \\
              & + InfoMap flow k3                & 27.12 & 91.94  & 11.54 & 1.24  \\
              & + METIS k3                       & 33.76 & 100.13 & 17.43 & 1.69  \\
    \midrule
    \multirow{3}{*}{MAPPO}
              & No partition (baseline)          & 29.69 & 94.54  & 13.66 & 1.42  \\
              & + InfoMap flow k3                & 25.27 & 90.20  & 9.63  & 1.12  \\
              & + METIS k3                       & 26.78 & 91.74  & 10.32 & 1.24  \\
    \midrule
    \multirow{3}{*}{IDQN}
              & No partition (baseline)$\dagger$ & 23.69 & 88.63  & 7.40  & 1.03  \\
              & + InfoMap flow k3                & 27.68 & 92.55  & 10.50 & 1.32  \\
              & + METIS k3                       & 25.26 & 90.29  & 8.73  & 1.18  \\
    \midrule
    \multirow{3}{*}{CoLight (seed 7)}
              & No partition (baseline)          & 33.58 & 98.64  & 16.20 & 1.77  \\
              & + InfoMap flow k3                & 46.81 & 112.22 & 28.06 & 2.69  \\
              & + METIS k3                       & 34.14 & 99.02  & 15.92 & 1.82  \\
    \bottomrule
\end{tabular}
\end{table}

\noindent\textbf{Notes.}
IDQN baseline crashed at episode 1340; its metrics are not comparable to
stable runs and IDQN is excluded from main-text analysis.
CoLight+flow k3 diverged (wait rises to 40--121\,s/veh after ep\,700);
root cause: $|M_k| \approx 2$--3 agents per module is insufficient for
CoLight's attention mechanism.
CoLight+METIS k3 matches the no-partition baseline within noise.
MAPPO+flow achieves the best wait time (9.63\,s/veh) across all valid conditions.

\subsection{Ingolstadt21 (I21, 21 intersections, seed\,0)}
\label{app:results:i21}

\begin{table}[H]
\centering
\caption{Ingolstadt21 MAPPO and IPPO --- all runs complete (1500\,ep).
    IPPO baseline crashed at ep\,1200 (marked $\dagger$).}
\label{tab:i21_mappo_ippo}
\small
\begin{tabular}{lllrrrr}
\toprule
Algorithm & Condition & $k$ & Delay & Trip & Wait & Queue \\
\midrule
\multirow{5}{*}{MAPPO}
& Global (no partition) & --- &  566.43 & 659.06 & 484.19 &  9.59 \\
& + InfoMap flow        & 4   &  191.25 & 326.29 & \textbf{140.96} &  4.44 \\
& + InfoMap flow        & 5   &  262.48 & 328.83 & 160.51 &  4.24 \\
& + METIS               & 4   &  607.96 & 547.39 & 401.17 &  7.88 \\
& + METIS               & 5   &  746.90 & 828.65 & 640.82 & 11.62 \\
\midrule
\multirow{5}{*}{IPPO}
& Local (no partition)$\dagger$ & --- & 908.49 & 870.11 & 733.19 & 11.28 \\
& + InfoMap flow        & 4   &  619.43 & 725.58 & 564.96 &  8.05 \\
& + InfoMap flow        & 5   &  253.14 & 382.59 & \textbf{202.71} &  4.89 \\
& + METIS               & 4   & 1048.57 & 962.23 & 846.14 & 12.14 \\
& + METIS               & 5   &  694.80 & 670.52 & 526.10 &  7.90 \\
\bottomrule
\end{tabular}
\end{table}

\begin{table}[H]
    \centering
    \caption{Ingolstadt21 CoLight --- all runs \textbf{incomplete} (target:
        1500\,ep). Episode counts indicate training progress at logging time.
        Two CoLight variants: Variant~A (PartitionedAttentionStack, shared
        network) and Variant~B (CoLightRegional, separate per-region learners).
        Seed 7 for Variant~B; seed varies for Variant~A.}
    \label{tab:i21_colight}
    \small
    \begin{tabular}{lllrrrrrr}
        \toprule
        Variant & Condition                 & $k$ & Delay   & Trip   & Wait            & Queue & Ep   \\
        \midrule
        \multirow{3}{*}{A (Partitioned)}
                & Fixed topology (baseline) & --- & 731.74  & 734.21 & 591.26          & 9.36  & 800  \\
                & + InfoMap flow            & --- & 535.18  & 526.03 & \textbf{365.26} & 7.46  & 620  \\
                & + METIS                   & --- & 953.24  & 832.20 & 702.45          & 13.01 & 630  \\
        \midrule
        \multirow{3}{*}{B (Regional)}
                & No partition (baseline)   & --- & 1031.66 & 870.62 & 747.70          & 9.91  & 1050 \\
                & + InfoMap flow            & 5   & 902.20  & 746.84 & \textbf{629.34} & 10.37 & 910  \\
                & + METIS                   & 5   & 944.03  & 784.10 & 656.28          & 9.43  & 820  \\
        \bottomrule
    \end{tabular}
\end{table}

\noindent\textbf{Notes.}
MAPPO+flow\,k4 (wait\,=\,140.96) outperforms the global no-partition MAPPO
(484.19) by $3.43\times$, and outperforms METIS\,k4 (401.17) by $2.85\times$.
IPPO+flow\,k5 (202.71) outperforms IPPO+METIS\,k5 (526.10) by $2.6\times$.
MAPPO flow benefits from coarser partitions ($k$\,=\,4); IPPO benefits from
finer partitions ($k$\,=\,5) --- consistent with the different roles of the
partition at each coupling level.
CoLight Variant~A with flow partition shows 365.26 wait at ep\,620, the
best CoLight result on I21; this run should be completed to 1500\,ep.
CoLight Variant~B (separate learners) is stable on I21 ($|M_k| \approx 4$
agents/module) unlike on C8 ($|M_k| \approx 2$); flow beats no-partition but
has not reached convergence.

\paragraph{Open runs (as of 2026-06-30).}
\begin{itemize}
    \item CoLight Variant~A flow (I21) --- \textbf{resume to 1500\,ep}
    \item CoLight Variant~A METIS (I21) --- resume to 1500\,ep
    \item MAPPO/IPPO no-partition baseline (I21 -- crashed runs)
    --- rerun with fixed seed
    \item Multi-seed expansion ($\geq$5 seeds): MAPPO+flow\,k4 on I21
    (primary condition for rliable evaluation)
\end{itemize}


\end{document}
